\documentclass[12pt]{amsart}
\pagestyle{empty}
\usepackage[hmargin=1in,vmargin=1in]{geometry}
\usepackage{url,  mathrsfs}
\usepackage{hyperref}
\usepackage{fancyhdr} 
\usepackage[T1]{fontenc}


\usepackage{listings,xcolor}

\lstset{language=Mathematica}
\lstset{basicstyle={\sffamily\footnotesize},
  numbers=left,
  numberstyle=\tiny\color{gray},
  numbersep=5pt,
  breaklines=true,
  captionpos={t},
  frame={lines},
  rulecolor=\color{black},
  framerule=0.5pt,
  columns=flexible,
  tabsize=2
}


\usepackage{url,  mathrsfs}
\usepackage{hyperref}
\usepackage{amsmath}
\usepackage{graphicx,  epstopdf}
\usepackage{subfig}
\usepackage{color}
\usepackage{bbm}
\usepackage{subfloat}
\usepackage[draft]{fixme}
\usepackage{bm} 
\usepackage{bbm} 
\usepackage{epigraph}
\usepackage{endnotes}
\usepackage{ifpdf}
\usepackage{array}

\usepackage{multicol}
\usepackage{multirow}
\usepackage{graphicx}
\usepackage{tikz}

\newcommand{\A}{\mathbb{A}}
\newcommand{\Q}{\mathbb{Q}}
\newcommand{\N}{\mathbb{N}}
\newcommand{\Z}{\mathbb{Z}}
\newcommand{\R}{\mathbb{R}}
\newcommand{\F}{\mathbb{F}}
\newcommand{\C}{\mathbb{C}}
\renewcommand{\P}{\mathbb{P}}
\newcommand{\I}{\mathcal{I}}
\newcommand{\V}{\mathcal{V}}
\newcommand{\cH}{\mathcal{H}}
\newcommand{\ol}{\overline}
\newcommand{\J}{\mathcal{J}}
\newcommand{\X}{\mathcal{X}}
\theoremstyle{definition}
\newtheorem{prob}{Problem}
\newtheorem{extraProb}{Extra Problem}

\DeclareMathOperator{\conv}{conv}
\DeclareMathOperator{\cone}{cone}
\DeclareMathOperator{\ext}{ext}




\begin{document}
 \thispagestyle{empty}
\begin{center} {\large \textbf{Math 380 -- Homework  6}} \smallskip \\ 
Due on Thursday, May 21\ \\ \ \\
 \end{center}

You are welcome to talk with other students in the class about problems but should write up solutions on your own. 
Solutions can be handwritten or typed but need to be legible and submitted via Gradescope by the end of the day 
on Thursday. You should justify all your answers in order to receive full credit.  
\bigskip 

\begin{flushleft}

{\bf Good practice problems} (not to be turned in): \\

CLO Section 4.1, Exercises 1, 7 \\
CLO Section 4.2, Exercises 1, 4, 5, 6, 7, 14, 15, 17 \\
\ \ \\
\bigskip 




\begin{prob} CLO \S4.1, Exercise 2 \smallskip \\ 
\begin{enumerate}
\item[(a)] Consider $J = \langle x^2+y^2-1,y-1\rangle$. Then $\mathbf V(J) = \{ (0, 1) \}$, so clearly $x \in \mathbf I(\mathbf V(J))$. However, we can compute that $G=\langle y-1,x^2\rangle$ is a Gr\"obner basis for $J$ (under lexicographical $x > y$ order), so since $\overline{x}^G = x \neq 0$ we know $x \not\in J$. 

\begin{lstlisting}[language=Mathematica]
GroebnerBasis[{x^2 + y^2 - 1, y - 1}, {x, y}]
\end{lstlisting}
\item[(b)] And \emph{What is the smallest $m$ for which $f^m\in J$?} $m=2$: $x^2 \in J = \langle y-1,x^2\rangle$.
\end{enumerate}
\end{prob}

\bigskip 


\newpage{}
\begin{prob}  CLO \S4.2, Exercise 16 \smallskip \\
Let $I \subseteq k[x_1,\dots,x_n]$ be an ideal with a Gr\"obner basis $G = \{g_1,\dots,g_s\}$ such that for all $i$, $\text{LT}(g_i)$ is square free in that $\langle \text{LT}(g_i)\rangle = \sqrt{\langle\text{LT}(g_i\rangle}$.
\begin{enumerate}
    \item[(a)] If $f \in \sqrt{I}$, then $\mathrm{LT}(f)$ is divisible by $\text{LT}(g_i)$ for some $i$. Let $m \in \mathbb N_{>0}$ be such that $f^m \in I$. Since $G$ is a Gr\"obner basis, some $\mathrm{LT}(g_i)$ must divide $\text{LT}(f^m)$. Then since $\text{LT}(a\cdot b) = \text{LT}(a)\cdot \text{LT}(b)$ by our proof in some previous homework, $\text{LT}(f^m) = \text{LT}(f)^m$. Clearly since $\text{LT}(g_i)$ is square free and divides $\text{LT}(f)^m$, $\text{LT}(g_i)$ must also divide $\text{LT}(f)$ (more details below).
    
    We can drop the scalar factors and consider the monomial part of $\text{LM}(f) = x^\alpha$ and $\text{LM}(g_i) = x^\beta$, and since $\text{LT}(g_i)$ is square free we know $\beta_j \in \{0,1\}$ for all $j$. Since $\text{LT}(g_i)$ divides $\text{LT}(f)^m$, we know that for some $h \in k[x_1,\dots,x_n]$, we have $h\cdot x^{\beta} = x^{m \cdot \alpha}$. Here, since we are dealing with monomials, $h(x)$ must be a monomial too, and we can write $x^{\beta+\gamma} = x^{m \cdot \alpha}$ for some $\gamma \in \mathbb Z^n_{\ge 0}$. The fact that $\beta + \gamma = m\cdot \alpha$ means that $\beta_j\le m \cdot \alpha_j$ for all $j$, and this can only be true when $\alpha_j > 0$ for all $j$ where $\beta_j =1$. Thus, $\beta_j \le \alpha_j$, and so $x^\beta \mid x^\alpha$, and $\text{LT}(g_i)$ divides $\text{LT}(f)$.
    
    
    \item[(b)] $I$ is radical. First, note that since $G \subseteq I\subseteq \sqrt I$, $\langle G \rangle \subseteq \sqrt I$ and $\langle \text{LT}(G)\rangle \subseteq \langle \text{LT}(\sqrt I)\rangle$. By part (a) above, for any $f \in \sqrt I$ we have that some $\text{LT}(g_i)$ divides $\text{LT}(f)$, so $\langle \text{LT}(\sqrt I)\rangle \subseteq \langle \text{LT}(G)\rangle $. Thus, $G$ is also a Gr\"obner basis for $\sqrt I$: dividing any $h \in \sqrt I$ by $G$ will necessarily give 0 remainder, so $\sqrt{I} \subseteq \langle G \rangle$, and in fact $\sqrt I = \langle G\rangle$. Finally, since $\langle G\rangle = I = \sqrt{I}$, by part (b) of Exercise 4 of \S4.2, $I$ is radical.
    \item[(c)] And \emph{Given an example of a radical ideal $I$ for which there is \emph{not} a Gr\"obner basis with square-free leading terms. 
(\emph{Hint}: One can take $I\subseteq k[x]$.)} 

    Take $I = \langle(x-2)(x-3)\rangle \subseteq k[x]$. We can show that $I=\mathbf I(\{2,3\})$. Any polynomial $f \in \mathbf I(\{2,3\})$ must vanish on $2$ and $3$, so by the argument with the division algorithm it must be divisible both by $x-2$ and $x-3$. Then, $f$ must be divisible by $(x-2)(x-3)$ (polynomial rings are UFDs and $(x-2)$ and $(x-3)$ are relatively prime?). Thus, we have $\mathbf I(\{2,3\})\subseteq I$, and the opposite inclusion is obvious since $(x-2)(x-3)$ vanishes on $\{2,3\}$. Since any ideal of a variety is radical, $I$ is radical with $I = \sqrt I$. However, since $I = \langle x^2-5x+6\rangle$ and $\text{LT}(x^2-5x+6)=x^2$, every element $f \in I$ will have $\text{LT}(f)$ divisible by $x^2$, as $f = h\cdot (x^2-5x+6)$ for some $h \in k[x]$, and $\text{LT}(f) = \text{LT}(h)\cdot x^2$. Thus, there will be no Gr\"obner basis for $I$ with generators having square-free leading terms.
\end{enumerate}
\ \\ \ \\
You may use the results of \S4.2, Exercise 4 without proof.
\end{prob}

\bigskip 


\begin{prob}  By the end of Friday, May 22, finalize a choice for your final project by emailing me (vinzant@uw.edu) with a topic, reference, and list of your group members. \end{prob}


\end{flushleft}
\end{document}